%% 350 words or less. Present tense, third person. No page number.

% In split inference, the large language model (LLM) is partitioned into blocks, with the cloud running the middle layers.

Split inference partitions a large language model (LLM) across client and cloud: the client runs the early and final layers, and the cloud runs the middle layers between them. One cloud deployment then serves many clients, which lowers the cost of inference, but every client sends the hidden states of its prompts to the cloud. Those hidden states need protection, because prior work recovers the prompt from them. Each such inverter is trained on the hidden states of one known model, so a target model whose weights the attacker does not hold is out of reach.

This dissertation reaches the target model through a shared space. Following Jha et al.~\todo{cite Jha et al.}, who show that hidden states of different models align in one space, a translation hub joins Qwen3.5-0.8B, OLMo-2-1B, and SmolLM2-1.7B: model-specific adapters and a shared trunk map each model's hidden states into the hub, and model-specific projectors condition each frozen model to reconstruct the prompt. Training the hub end to end on reconstruction raises cross-model recovery to the fully supervised level, word-F1 of 0.95 on question prompts and 0.98 on legal text, where any model's hidden states reconstruct the prompt through any other model's decoder.

The hub then becomes the inversion target for a model the attacker does not hold. The attacker knows the target model's tokenizer and vocabulary, holds none of its weights, issues no queries, and receives a small leak: hidden-state rows at single token positions, each paired with the token that produced it. A lightweight affine map, trained on the leaked rows, carries the target model's hidden states into the hub, and the frozen decoders reconstruct the prompt.

The leak is small, and matching makes it sufficient. The leaked rows anchor the map; unsupervised matching on unpaired hidden states extends it. With 90 leaked rows, nine from each of ten prompts, matching raises word-F1 on Pythia-1B from 0.35 to 0.77 on question prompts and from 0.10 to 0.32 on legal text. With 864 rows it raises word-F1 from 0.76 to 0.81 and from 0.55 to 0.63, so the gain from matching shrinks as the leak grows.

The remaining work asks how much matching contributes and where its gain ends, under three tests: leak budget, corpus domain, and target model. The answers measure how little a split-inference system can leak before its hidden states reveal the prompt.

%%------------------- note
% Hub - balancing of supervised - unsupervised; reconstruct

% 1. Attack - privacy focusing; front load impact on how, privacy is important. Anthropic scoop - navius. Information protection at local users. news -->reserach.
%     1b. why split inference, red teaming

% 2. Based on platonic representation;
%     2a. few pieces of collection --> surprising results; statement of LLM has common --> leading to attack


% 3. Beyond the paper

% impact
% red teaming - split inference

% current results

% follow up -- build, more structure, sequence, architecture, chronically

% 1 example red team specific methods -- NOIR like; combination of strategy

% ground floor - generalize into how LLM representation used, applied hidden state

% ------
% writing + (10) bibliography

% main focus -- method -- results